\documentclass[11pt]{article}
\usepackage[utf8]{inputenc}
\usepackage{amsmath, amssymb, amsthm}
\usepackage[margin=1in]{geometry}
\usepackage{hyperref}

\newtheorem{theorem}{Theorem}
\newtheorem{definition}{Definition}
\newtheorem{lemma}[theorem]{Lemma}
\newtheorem{proposition}[theorem]{Proposition}

\title{Integration (Riemann)}
\author{math-foundations concept 15}
\date{}

\begin{document}
\maketitle

\subsection*{First, an example}
Before the abstract Darboux machinery, here is the punchline on a concrete case.
Take $f(x) = x^2$ on $[0,1]$. Slice $[0,1]$ into $n=4$ equal subintervals of width
$\Delta x = 1/4$ and form the right-endpoint Riemann sum
\[ R_4 = \sum_{i=1}^{4} \left(\tfrac{i}{4}\right)^2 \cdot \tfrac{1}{4}
       = \tfrac{1}{4}\!\left(\tfrac{1}{16} + \tfrac{4}{16} + \tfrac{9}{16} + \tfrac{16}{16}\right)
       = \tfrac{30}{64} = 0.46875. \]
\emph{Read aloud:} R-sub-4 equals the sum from i equals 1 to 4 of i-over-4 squared, times one-quarter; this works out to thirty over sixty-four, about 0.469.
Refining to $n=100$ gives $\approx 0.3383$; the limit is $\int_0^1 x^2\,dx = \tfrac{1}{3}$,
which the antiderivative $F(x) = x^3/3$ recovers instantly: $F(1) - F(0) = 1/3$.
Everything below explains \emph{why} that limit exists and \emph{why} the antiderivative
shortcut works.

\section{Motivation}
The Riemann integral assigns to a bounded function $f : [a,b] \to \mathbb{R}$ a single
number $\int_a^b f$ that represents the signed area between the graph of $f$ and the
$x$-axis. The construction is purely metric: chop $[a,b]$ into finitely many subintervals,
approximate $f$ on each subinterval by its supremum (over\-estimate) and infimum
(under\-estimate), and demand that the two estimates converge as the chopping becomes
finer.

\section{Partitions and Darboux Sums}

\begin{definition}[Partition]
A \emph{partition} of $[a,b]$ is a finite ordered set
$P = \{x_0, x_1, \dots, x_n\}$ with
\[ a = x_0 < x_1 < \cdots < x_n = b. \]
\emph{Read aloud:} a equals x-zero, less than x-one, less than dot-dot-dot, less than x-n, equals b.
Set $\Delta x_i = x_i - x_{i-1}$ and define the \emph{mesh} $\|P\| = \max_i \Delta x_i$.
A partition $P'$ \emph{refines} $P$ if $P \subseteq P'$.
\end{definition}

\begin{definition}[Upper and lower sums]
Let $f : [a,b] \to \mathbb{R}$ be bounded. For a partition $P$ as above, set
\[ M_i = \sup_{x \in [x_{i-1}, x_i]} f(x), \qquad m_i = \inf_{x \in [x_{i-1}, x_i]} f(x). \]
\emph{Read aloud:} big-M-sub-i is the supremum of f of x for x in the i-th subinterval; little-m-sub-i is the infimum on the same subinterval.
The \emph{upper Darboux sum} and \emph{lower Darboux sum} are
\[ U(f, P) = \sum_{i=1}^{n} M_i\, \Delta x_i, \qquad L(f, P) = \sum_{i=1}^{n} m_i\, \Delta x_i. \]
\emph{Read aloud:} U of f comma P equals the sum from i equals 1 to n of big-M-sub-i times delta-x-sub-i; L of f comma P is the analogous sum with little-m-sub-i.
\end{definition}

\begin{lemma}\label{lem:refine}
If $P'$ refines $P$, then $L(f, P) \le L(f, P') \le U(f, P') \le U(f, P)$.
Consequently $L(f, P_1) \le U(f, P_2)$ for any two partitions $P_1, P_2$.
\end{lemma}

\begin{definition}[Darboux integrability]
Define the \emph{upper} and \emph{lower integrals}
\[ \overline{\int_a^b} f = \inf_P U(f, P), \qquad \underline{\int_a^b} f = \sup_P L(f, P). \]
\emph{Read aloud:} upper integral from a to b of f equals the infimum over partitions P of U of f, P; lower integral equals the supremum over P of L of f, P.
We say $f$ is \emph{Riemann (Darboux) integrable} on $[a,b]$ if the two coincide; the
common value is denoted $\int_a^b f(x)\,dx$.
\end{definition}

\begin{proposition}[Cauchy criterion]\label{prop:cauchy}
A bounded $f$ is integrable on $[a,b]$ iff for every $\varepsilon > 0$ there exists a
partition $P$ with $U(f, P) - L(f, P) < \varepsilon$.
\end{proposition}

\begin{proof}
($\Rightarrow$) Suppose $\overline{\int} f = \underline{\int} f =: I$. By the
definitions of $\inf$ and $\sup$, pick $P_1, P_2$ with $U(f, P_1) < I + \varepsilon/2$
and $L(f, P_2) > I - \varepsilon/2$. The common refinement $P = P_1 \cup P_2$ satisfies
$U(f, P) \le U(f, P_1)$ and $L(f, P) \ge L(f, P_2)$, so $U(f,P) - L(f,P) < \varepsilon$.

($\Leftarrow$) For every $\varepsilon$ we have
$0 \le \overline{\int} f - \underline{\int} f \le U(f, P) - L(f, P) < \varepsilon$.
Hence the difference is $0$.
\end{proof}

\section{Continuity Implies Integrability}

\begin{theorem}\label{thm:main}
Every continuous function $f : [a, b] \to \mathbb{R}$ is Riemann integrable.
\end{theorem}

\begin{proof}
Because $[a,b]$ is compact, $f$ is bounded and uniformly continuous (Heine--Cantor).
Fix $\varepsilon > 0$. By uniform continuity, choose $\delta > 0$ such that
\[ x, y \in [a,b], \; |x - y| < \delta \implies |f(x) - f(y)| < \frac{\varepsilon}{b - a}. \]
\emph{Read aloud:} for all x, y in the closed interval a to b, if the absolute value of x minus y is less than delta, then the absolute value of f of x minus f of y is less than epsilon over b minus a.
Take any partition $P = \{x_0, \dots, x_n\}$ with mesh $\|P\| < \delta$.

On each subinterval $[x_{i-1}, x_i]$, the function $f$ attains its sup $M_i$ and inf
$m_i$ (continuous on a compact set), say at points $\xi_i, \eta_i \in [x_{i-1}, x_i]$.
Since $|\xi_i - \eta_i| \le \Delta x_i < \delta$,
\[ M_i - m_i = f(\xi_i) - f(\eta_i) < \frac{\varepsilon}{b - a}. \]
\emph{Read aloud:} big-M-sub-i minus little-m-sub-i equals f of xi-sub-i minus f of eta-sub-i, which is less than epsilon over b minus a.

Summing over $i$,
\[ U(f, P) - L(f, P)
   = \sum_{i=1}^{n} (M_i - m_i)\, \Delta x_i
   < \frac{\varepsilon}{b - a} \sum_{i=1}^{n} \Delta x_i
   = \frac{\varepsilon}{b - a} \cdot (b - a)
   = \varepsilon. \]
\emph{Read aloud:} U of f, P minus L of f, P equals the sum from i equals 1 to n of M-sub-i minus m-sub-i, times delta-x-sub-i, which is less than epsilon over b minus a, times the total length b minus a, equals epsilon.

By Proposition~\ref{prop:cauchy}, $f$ is integrable.
\end{proof}

\section{The Fundamental Theorem of Calculus}

\begin{theorem}[FTC, Part 1]
Let $f : [a, b] \to \mathbb{R}$ be continuous, and define $F : [a, b] \to \mathbb{R}$ by
$F(x) = \int_a^x f(t)\,dt$. Then $F$ is differentiable on $(a, b)$ and $F'(x) = f(x)$.
\end{theorem}

\begin{proof}[Sketch]
For $h \neq 0$,
$\frac{F(x+h) - F(x)}{h} = \frac{1}{h} \int_x^{x+h} f(t)\,dt$.
Bound the integrand by $f(x) \pm \varepsilon$ using continuity of $f$ at $x$, and let
$h \to 0$.
\end{proof}

\begin{theorem}[FTC, Part 2]
If $f$ is continuous on $[a,b]$ and $G$ is any antiderivative of $f$, then
$\int_a^b f(t)\,dt = G(b) - G(a)$.
\end{theorem}

\begin{proof}[Sketch]
Both $F$ from Part 1 and $G$ are antiderivatives of $f$, so they differ by a constant:
$G(x) = F(x) + C$. Then $G(b) - G(a) = F(b) - F(a) = \int_a^b f - 0$.
\end{proof}

\section{Worked Example}
For $f(x) = x^2$ on $[0,1]$ with equal partition $x_i = i/n$,
\[ U(f, P_n) = \sum_{i=1}^{n} \frac{i^2}{n^2}\cdot\frac{1}{n}
              = \frac{(n+1)(2n+1)}{6n^2} \xrightarrow[n\to\infty]{} \frac{1}{3}. \]
\emph{Read aloud:} U of f, P-sub-n equals the sum from i equals 1 to n of i-squared over n-squared, times one over n; this simplifies to (n plus 1)(2n plus 1) over 6 n-squared, which tends to one-third as n goes to infinity.
By FTC2 the antiderivative $x^3/3$ also yields $1/3 - 0 = 1/3$.

\end{document}
